\section{Hạn chế và hướng phát triển}

\subsection{Hạn chế hiện tại}

Hạn chế đầu tiên đến từ chính Dark Channel Prior. Phương pháp này hoạt động tốt trên nhiều ảnh ngoài trời, nhưng kém ổn định ở các vùng trời sáng, bề mặt trắng lớn hoặc các vùng không chứa pixel tối theo đúng giả thiết của prior. Trong các trường hợp đó, transmission có thể bị ước lượng lệch và ảnh sau khử haze dễ mất tự nhiên.

Hạn chế thứ hai nằm ở hiện thực plugin. Toàn bộ pipeline hiện tại chạy trên CPU và thực hiện nhiều phép lọc cục bộ trên toàn ảnh. Cơ chế cache giúp giảm đáng kể tính toán lặp, nhưng khi áp dụng lên ảnh lớn, full render vẫn có thể chậm nếu so với kỳ vọng của một công cụ chỉnh sửa ảnh thời gian thực.

Thứ ba, adaptive omega hiện mới dựa trên một mô hình hồi quy rất gọn với ba đặc trưng: dark channel, contrast và saturation. Cách tiếp cận này có ưu điểm là dễ tích hợp, chi phí thấp và giải thích được, nhưng cũng giới hạn khả năng mô tả các trường hợp phức tạp hơn của ảnh haze thực tế. Các hệ số hồi quy đang được cố định theo mô hình đã huấn luyện trong notebook, nên chất lượng của adaptive mode vẫn phụ thuộc vào mức độ khớp giữa ảnh thực tế và dữ liệu đã dùng để học.

Cuối cùng, phần đánh giá trong plugin hiện chủ yếu dựa trên quan sát trực quan và thử nghiệm tương tác. Điều này phù hợp với mục tiêu ứng dụng, nhưng vẫn chưa thay thế được một quy trình benchmark hệ thống hơn trên tập ảnh đa dạng có ground truth hoặc ít nhất có giao thức so sánh nhất quán giữa nhiều cấu hình tham số.

\subsection{Hướng phát triển}

Hướng phát triển trực tiếp nhất là học lại hoặc hiệu chỉnh các hệ số của \(\omega_{\mathrm{map}}\) trên một tập dữ liệu phong phú hơn, bao phủ nhiều tình huống như vùng trời lớn, cảnh ngược sáng, ảnh đô thị có nhiều bề mặt phản xạ và ảnh có màu sắc ít bão hòa. Điều này có thể giúp nhánh adaptive ổn định hơn khi chuyển từ môi trường notebook sang ảnh thực tế trong plugin.

Một hướng khác là mở rộng bộ đặc trưng đầu vào cho adaptive omega. Ngoài dark channel, contrast và saturation, có thể cân nhắc thêm đặc trưng kết cấu, độ sáng cục bộ hoặc các tín hiệu gợi ý vùng trời. Mục tiêu là tăng khả năng phân biệt giữa vùng cần khử haze mạnh và vùng chỉ nên xử lý nhẹ, nhưng vẫn giữ mô hình đủ gọn để chạy hiệu quả trong plugin.

Về hiệu năng, có thể tối ưu thêm các phép lọc cửa sổ, song song hóa tốt hơn các vòng lặp theo pixel hoặc cân nhắc khai thác GPU nếu môi trường Paint.NET và API cho phép. Đây là hướng quan trọng nếu muốn biến plugin từ một nguyên mẫu học thuật thành công cụ có tính thực dụng cao hơn trên ảnh độ phân giải lớn.

Ở phía giao diện, plugin có thể được mở rộng bằng các tùy chọn như bật/tắt guided refinement, thay đổi kích thước patch, hiển thị so sánh song song Fixed/Adaptive hoặc cho phép xem trước transmission map. Các tùy chọn này không làm thay đổi bản chất thuật toán, nhưng giúp người dùng hiểu rõ hơn tác động của từng thành phần trong pipeline.

Xa hơn nữa, nếu mục tiêu chuyển sang chất lượng khôi phục tối đa thay vì tính gọn nhẹ của hệ thống, có thể cân nhắc các mô hình mạnh hơn như phương pháp học sâu cho dehazing. Tuy nhiên, đối với phạm vi đồ án hiện tại, hướng phát triển hợp lý hơn vẫn là tiếp tục cải tiến DCP theo hướng có thể giải thích được, dễ triển khai và phù hợp với môi trường plugin thực tế.
